\section{Wnioski}
\label{sec:wnioski}

%<*wnioski>
\begin{enumerate}
  \item \textbf{Nie ma jednej najszybszej bazy.} Valkey wygrywa w~operacjach
  po kluczu, w~licznikach i~w~rezerwacji miejsc przy dużej współbieżności.
  Bez przygotowanych przez aplikację struktur dostępu jest jednak najwolniejszy
  ze wszystkich: historia klienta (R4) bez indeksów trwa w~nim
  \val{valkey}{10000000}{off}{R4}\,ms, a~w~PostgreSQL \val{postgres}{10000000}{off}{R4}\,ms. PostgreSQL jest najbardziej
  przewidywalny: bez indeksów najszybszy z~baz dyskowych, z~indeksami
  w~czołówce wszystkich scenariuszy.
  \item \textbf{Indeksy decydują bardziej niż silnik.} Ten sam scenariusz
  w~tej samej bazie przyspiesza po dodaniu indeksu nawet o~kilka rzędów
  wielkości (rys.~\ref{fig:speedup}). Po zaindeksowaniu różnice między
  silnikami w~zapytaniach selektywnych są zwykle mniejsze niż rząd wielkości.
  Duże pozostają w~agregacjach po całej tabeli (R6, J2). Koszt to wolniejsze wstawianie hurtowe
  i~dodatkowe miejsce.
  \item \textbf{Model danych zmienia złożoność operacji.} Osadzenie biletów
  w~dokumencie rezerwacji daje w~MongoDB odczyt i~anulowanie rezerwacji
  w~czasie stałym bez indeksów drugorzędnych. Baza relacyjna potrzebuje do
  tego indeksu na kluczu obcym.
  \item \textbf{Współbieżność to kwestia strategii blokowania.} Przy 1000
  klientów strategia (warunkowy UPDATE, \texttt{FOR UPDATE NOWAIT}, skrypt
  Lua, WATCH) zmienia przepustowość tego samego silnika nawet o~trzy
  rzędy wielkości (Valkey: \texttt{lua} a~\texttt{watch}). Szybkie odrzucanie konfliktów (\texttt{NOWAIT}, operacja
  atomowa) skaluje się lepiej niż kolejkowanie na blokadach.
  \item \textbf{Bezpieczeństwo bez kosztu, z~jednym wyjątkiem.} Konta
  o~minimalnych uprawnieniach i~TLS nie zmieniają czasu operacji CRUD ponad
  zmienność pomiaru. Wyjątkiem jest Valkey przy 100 klientach, gdzie TLS
  obniża przepustowość z~\seccval{valkey}{base}{100} do
  \seccval{valkey}{tls}{100}\,rez./s (rozdz.~\ref{sec:bezpieczenstwo}).
  \item \textbf{H1 potwierdzona dla Valkey, częściowo dla MongoDB, H2
  potwierdzona częściowo} (rozdz.~\ref{sec:hipotezy}).
  \item \textbf{Praktyczna rekomendacja dla systemu biletowego:}
  rezerwacje i~płatności w~PostgreSQL (spójność, złożone zapytania), blokady
  miejsc z~TTL i~liczniki dostępności w~Valkey (operacja atomowa w~skrypcie
  Lua). Katalog wydarzeń może być w~MongoDB, jeśli jego struktura często się
  zmienia.
\end{enumerate}
%</wnioski>

\subsection{Zagrożenia dla trafności wyników}
\label{sec:zagrozenia}
%<*zagrozenia>

\begin{itemize}
  \item Wszystkie bazy działały na jednym komputerze w~maszynie wirtualnej
  Docker Desktop (8\,GB RAM, Apple M5). Wyniki pokazują relacje między
  silnikami, a~nie przepustowość produkcyjną. Generator obciążenia działał na
  tej samej maszynie co bazy.
  \item Poziomy trwałości nie są równe: Valkey bez RDB i~AOF nie czeka na dysk,
  MongoDB czeka na dziennik (\texttt{majority}), bazy relacyjne na
  \texttt{fsync} WAL i~redo. Zapisy w~Valkey są przez to tańsze z~założenia.
  \item Każda baza działała jako pojedynczy węzeł, a~główna macierz bez TLS.
  Wpływ TLS i~ról zmierzyliśmy osobno (rozdz.~\ref{sec:bezpieczenstwo}),
  wpływu replikacji synchronicznej nie badaliśmy.
  \item W~każdym wariancie mierzyliśmy 5 powtórzeń po jednym przebiegu
  rozgrzewkowym. Dla operacji trwających poniżej 1\,ms pojedyncze odchylenia
  (zimna pamięć podręczna, punkt kontrolny) przesuwają średnią, dlatego
  w~danych są też mediany i~p95.
  \item Indeksy w~Valkey są projektem autorów (zbiory \texttt{ix:*}), a~nie
  funkcją silnika. Inny projekt struktur dałby inne wyniki.
  \item Parametry konfiguracyjne baz zostały ustawione tak samo (1\,GB
  pamięci podręcznej), ale nie były strojone pod obciążenie. Strojenie
  (np. \texttt{innodb\_thread\_concurrency}, pooler połączeń dla PostgreSQL)
  zmieniłoby wyniki przy 1000 klientach.
\end{itemize}
%</zagrozenia>
